%% Question bank chapter file - generated by bank_tools.py build (do not edit while the bank is still being built:
%% it is regenerated from the per-test data in Question Bank\data). Every question is one bankq block:
%%   \begin{bankq}[id=..., source={...}, paper=.., qno=.., type=..., subtopic={...}, answer={...}, answer_src=key|solved,
%%                 working={...}, difficulty=1-10, diff_parts={K S I C F}, diff_note={...},
%%                 group=<passage id or empty>, also_in={duplicate copies in other tests}]
%% difficulty = min(10, K+S+I+C+F-1): K concepts 1-3, S steps 1-4, I insight 0-2, C computation 0-2, F format 0-1
%%   ...question LaTeX, options, figures (\bfig / \bfigi -> img/)...
%%   \end{bankq}
%% Passages shared by several questions are bankpassage blocks; their questions carry group=<passage id>.
\documentclass[10pt]{article}
\usepackage[a4paper,margin=1.8cm]{geometry}
\usepackage[T1]{fontenc}
\usepackage{amsmath,amssymb,mathtools}
\usepackage{graphicx}
\usepackage[export]{adjustbox}
\usepackage{xcolor,needspace,keyval}
\usepackage{enumitem}
\usepackage{booktabs,array,tabularx}
\usepackage[hidelinks]{hyperref}
\setlength{\parindent}{0pt}
\setlength{\parskip}{3pt}
\makeatletter
\define@key{bq}{id}{\def\bq@id{#1}}\define@key{bq}{source}{\def\bq@src{#1}}\define@key{bq}{paper}{\def\bq@paper{#1}}
\define@key{bq}{qno}{\def\bq@qno{#1}}\define@key{bq}{type}{\def\bq@type{#1}}\define@key{bq}{subtopic}{\def\bq@sub{#1}}
\define@key{bq}{answer}{\def\bq@ans{#1}}\define@key{bq}{answer_src}{\def\bq@asrc{#1}}\define@key{bq}{working}{\def\bq@work{#1}}
\define@key{bq}{group}{\def\bq@group{#1}}\define@key{bq}{also_in}{\def\bq@also{#1}}
\define@key{bq}{difficulty}{\def\bq@diff{#1}}\define@key{bq}{diff_parts}{\def\bq@dparts{#1}}\define@key{bq}{diff_note}{\def\bq@dnote{#1}}
\newenvironment{bankq}[1][]{\def\bq@work{}\def\bq@also{}\def\bq@group{}\def\bq@diff{?}\def\bq@dparts{}\def\bq@dnote{}\setkeys{bq}{#1}%
  \par\needspace{4\baselineskip}\medskip\noindent{\small\textbf{\bq@id}\hfill\color{gray}[\bq@src, Paper \bq@paper, Q\bq@qno]
  \bq@type\ $\cdot$ \bq@sub\ $\cdot$ \textbf{D\bq@diff}}\par\nopagebreak\relax}
  {\par{\small\color{teal}Answer: \textbf{\bq@ans}\ifx\bq@work\@empty\else\ (solved: \bq@work)\fi
  \ifx\bq@also\@empty\else\par\color{gray}Also in: \bq@also\fi}\par}
\define@key{bp}{id}{\def\bp@id{#1}}\define@key{bp}{source}{\def\bp@src{#1}}\define@key{bp}{questions}{\def\bp@qs{#1}}
\newenvironment{bankpassage}[1][]{\setkeys{bp}{#1}\par\needspace{4\baselineskip}\bigskip\noindent
  {\small\textbf{Passage \bp@id}\hfill\color{gray}[\bp@src; questions \bp@qs]}\par\nopagebreak\relax\itshape}{\par}
\makeatother
\newcommand{\bfig}[2]{\par\smallskip{\centering\includegraphics[width=#1pt,max width=0.62\linewidth,max height=0.32\textheight]{img/#2}\par}\smallskip}
\newcommand{\bfigi}[2]{\includegraphics[width=#1pt,max width=\linewidth,max height=0.32\textheight,valign=c]{img/#2}}
\newcommand{\bmissing}[1]{\textit{[figure missing in the source conversion: \detokenize{#1}]}}
\newenvironment{qoptions}{\begin{itemize}[leftmargin=1.8em,itemsep=1pt,topsep=2pt,label={}]}{\end{itemize}}
\newcommand{\opt}[2]{\item[(#1)] #2}
\providecommand{\fwstart}{}\providecommand{\fwend}{}\providecommand{\qfullsplit}{}
\newcolumntype{L}[1]{>{\raggedright\arraybackslash}p{#1}}
\newcolumntype{C}[1]{>{\centering\arraybackslash}p{#1}}
\newcolumntype{R}[1]{>{\raggedleft\arraybackslash}p{#1}}
\providecommand{\sectionhdr}[1]{}\providecommand{\subsectionhdr}[1]{}


\begin{document}
{\Large\bfseries Capacitance --- Question Bank}\par\smallskip
{\small Physics chapter P09. Questions from all converted test papers, grouped by source series. Answers: printed key where it matches the paper, otherwise solved (answer\_src=solved, with working). Exact repeats across papers appear once, with also\_in listing the other copies.}\par\medskip\hrule

\section*{Score-I papers}
\begin{bankq}[id=SI12-P1-PHY-S3-Q07, source={Score-I - Test 12 (09 Mar 2022)}, paper=1, qno=7, type=numerical, subtopic={Combinations}, answer={2}, answer_src=key, difficulty=2, diff_parts={K1 S1 I0 C0 F1}, diff_note={concepts: capacitors in steady state | steps: steady-state current through the resistors; potential across the 10 uF capacitor and charge}]
Find the charge in $\mu$C on $10\,\mu$F capacitor in steady state. If the value of charge is $x\,\mu$C, then find the value of $x/100$.
\bfig{92.3}{SI12_p1_phy_s3_q07_1.png}
\end{bankq}

\section*{TOAS papers}
\begin{bankq}[id=TOAS11-P1-PHY-S2-Q06, source={TOAS - Test 11 (27 Jun 2022)}, paper=1, qno=6, type=multi, subtopic={RC circuits}, answer={A,C}, answer_src=key, difficulty=3, diff_parts={K1 S2 I0 C0 F1}, diff_note={concepts: RC charging with an initially charged capacitor | steps: q(t) = epsilon C (1 - 2 e\^{}(-t/RC)); time when the charge is zero; final and maximum charge}]
A capacitor $C$ charged to a value $Q=\varepsilon C$ is connected to a battery of emf $\varepsilon$ as shown in the figure.
\bfig{115.6}{TOAS11_p1_phy_s2_q06_1.png}
\begin{qoptions}
\opt{A}{Charge on the capacitor is zero at time $RC\ln2$ after the switch S is closed.}
\opt{B}{Charge on the capacitor is zero very long time after the switch S is closed.}
\opt{C}{Maximum charge on the capacitor during the process of charging is $\varepsilon C$.}
\opt{D}{During the process of charging, the charging current first decreases and then increases exponentially.}
\end{qoptions}
\end{bankq}
\begin{bankq}[id=TOAS11-P1-PHY-S2-Q07, source={TOAS - Test 11 (27 Jun 2022)}, paper=1, qno=7, type=multi, subtopic={Charge sharing}, answer={A,B,C}, answer_src=key, difficulty=3, diff_parts={K1 S2 I0 C0 F1}, diff_note={concepts: charge redistribution among capacitors | steps: charge conservation on the connected plates; equal potentials; charge on each capacitor; direction of the flow}]
Two capacitors of 2 $\mu$F and 3 $\mu$F are charged to 150 volt and 120 volt respectively. The plates of a capacitor are connected as shown in the figure. A discharged capacitor of capacity 1.5 $\mu$F falls to the free ends of the wire. Then after the system comes in steady state.
\bfig{159.4}{TOAS11_p1_phy_s2_q07_1.png}
\begin{qoptions}
\opt{A}{Charge on the 1.5 $\mu$F capacitor will become 180 $\mu$C}
\opt{B}{Charge on the 2 $\mu$F capacitor will become 120 $\mu$C}
\opt{C}{Positive charge flows through A from left to right}
\opt{D}{Positive charge flows through A from right to left}
\end{qoptions}
\end{bankq}
\begin{bankq}[id=TOAS11-P1-PHY-S2-Q08, source={TOAS - Test 11 (27 Jun 2022)}, paper=1, qno=8, type=multi, subtopic={RC circuits}, answer={A,B,C}, answer_src=key, difficulty=3, diff_parts={K1 S2 I0 C0 F1}, diff_note={concepts: network with a capacitor | steps: current into the capacitor from dq/dt; potential relation; value of R and branch currents}]
The network shown in the figure is part of some bigger circuit. Charge on capacitor ($C=1\ \mu$F) at any time $t$ is $q=3(1-e^{-t})$ in coulombs. If at $t=0$ the potentials of point A and B is 20 V, then at $t=0$ match the following.
\bfig{157.9}{TOAS11_p1_phy_s2_q08_1.png}
\begin{qoptions}
\opt{A}{Current in branch AP is 3A}
\opt{B}{Current through R is 4 A}
\opt{C}{Value of R is 4.25 $\Omega$}
\opt{D}{Value of R is 1 $\Omega$}
\end{qoptions}
\end{bankq}
\begin{bankq}[id=TOAS11-P1-PHY-S2-Q11, source={TOAS - Test 11 (27 Jun 2022)}, paper=1, qno=11, type=multi, subtopic={Dielectrics}, answer={D}, answer_src=key, difficulty=3, diff_parts={K1 S2 I0 C0 F1}, diff_note={concepts: capacitor with a partial dielectric | steps: plates at the same potential so the field is the same in both parts; charge ratio; capacitance relation}]
A parallel plate capacitor has a dielectric slab of dielectric constant $K$ between its plates that cover 1/3 of the area of its plates, as shown in the figure. The total capacitance of the capacitor is $C$ while that of the portion with dielectric in between is $C_2$. When the capacitor is charged the plane area covered by the dielectric gets charge $Q_1$ and the rest of the area gets charge $Q_2$. The electric field in the dielectric is $E_1$ and that in the other portion is $E_2$. Choose the correct option(s), ignoring edge effects.
\bfig{91.3}{TOAS11_p1_phy_s2_q11_1.png}
\begin{qoptions}
\opt{A}{$\dfrac{E_1}{E_2}=K$}
\opt{B}{$\dfrac{E_1}{E_2}=\dfrac1K$}
\opt{C}{$\dfrac{Q_1}{Q_2}=\dfrac3K$}
\opt{D}{$\dfrac{C}{C_2}=\dfrac{2+K}{K}$}
\end{qoptions}
\end{bankq}
\begin{bankq}[id=TOAS11-P1-PHY-S2-Q12, source={TOAS - Test 11 (27 Jun 2022)}, paper=1, qno=12, type=multi, subtopic={Capacitors}, answer={B,C}, answer_src=key, difficulty=3, diff_parts={K1 S2 I0 C0 F1}, diff_note={concepts: capacitance of spherical and cylindrical and parallel-plate capacitors | steps: capacitance of the spherical capacitor; equal cylindrical capacitor length; equal parallel plate separation}]
A particular spherical capacitor consists of two concentric conducting spherical shells of radii $R$ and $2R$. Neglect edge effects in each.
\begin{qoptions}
\opt{A}{A coaxial cylindrical capacitor, made of two concentric cylindrical conductors of radii $R$ and $2R$, has the same capacitance as the spherical capacitor. Its length should be $4R/\ln2$}
\opt{B}{A coaxial cylindrical capacitor, made of two concentric cylindrical conductors of radii $R$ and $2R$, has the same capacitance as the spherical capacitor. Its length should be $4R\ln2$}
\opt{C}{A parallel-plate capacitor, made of two circular plates of radius $R$, has the same capacitance as the spherical capacitor. The plates are at a distance of $R/8$.}
\opt{D}{A parallel-plate capacitor, made of two circular plates of radius $R$, has the same capacitance as the spherical capacitor. The plates are at a distance of $R/4$.}
\end{qoptions}
\end{bankq}
\begin{bankq}[id=TOAS11-P2-PHY-S1-Q01, source={TOAS - Test 11 (27 Jun 2022)}, paper=2, qno=1, type=multi, subtopic={Dielectrics}, answer={B,C}, answer_src=key, difficulty=3, diff_parts={K1 S2 I0 C0 F1}, diff_note={concepts: capacitor with three dielectric slabs | steps: series and parallel groups; equivalent capacitance; fields in each region}]
An ideal parallel plate capacitor of area $A$ is filled with three dielectric slabs having dielectric constants $K_1=3.0$, $K_2=5.0$ and $K_3=2.0$ as shown in the figure. Choose the CORRECT option(s) :-
\bfig{103.0}{TOAS11_p2_phy_s1_q01_1.png}
\begin{qoptions}
\opt{A}{the equivalent capacitance of the system is $\dfrac{8\varepsilon_0A}{3d}$}
\opt{B}{the electric field in $K_1$ is more than the electric field in $K_2$.}
\opt{C}{the electric field on the left half of $K_3$ is less than the electric field on the right half of $K_3$.}
\opt{D}{the electric field on the left half of $K_3$ is more than the electric field on the right half of $K_3$.}
\end{qoptions}
\end{bankq}
\begin{bankq}[id=TOAS11-P2-PHY-S1-Q02, source={TOAS - Test 11 (27 Jun 2022)}, paper=2, qno=2, type=multi, subtopic={Charge sharing}, answer={A,D}, answer_src=key, difficulty=4, diff_parts={K1 S2 I0 C1 F1}, diff_note={concepts: charge sharing in a capacitor network | steps: initial charges; common potential before the battery; potential with the battery and the extra capacitor; charges and ground flow}]
Three capacitors of 2$\mu$F, 3$\mu$F and 5$\mu$F are individually charged with batteries of emf's 5V, 20V and 10V, respectively. After disconnecting from the voltage sources, these capacitors are connected as shown in the figure with their positive polarity plates connected to A and their negative polarity plates being earthed. Now, a battery of 20 V and an uncharged capacitor of 4 $\mu$F capacitance are connected to the junction A with a switch S, as shown. The switch is closed.
\bfig{163.9}{TOAS11_p2_phy_s1_q02_1.png}
\begin{qoptions}
\opt{A}{The potential of the junction A becomes $\dfrac{100}{7}$V}
\opt{B}{Final charge on the capacitor of lowest capacitance is 200 $\mu$C}
\opt{C}{The total charge flow to the ground after closing the switch is 200 $\mu$C}
\opt{D}{The charge flow to the ground after closing the switch is zero.}
\end{qoptions}
\end{bankq}
\begin{bankq}[id=TOAS11-P2-PHY-S1-Q05, source={TOAS - Test 11 (27 Jun 2022)}, paper=2, qno=5, type=multi, subtopic={Energy}, answer={A,C,D}, answer_src=key, difficulty=3, diff_parts={K1 S2 I0 C0 F1}, diff_note={concepts: energy and force of a capacitor | steps: heat dissipated while charging; force between plates; work when the plate separation is halved; dissipation}]
A parallel plate air capacitor has plate area $A$ and separation between the plates $d$. Switch S is closed to connect the capacitor to a cell of emf $V$.
\bfig{101.4}{TOAS11_p2_phy_s1_q05_1.png}
\begin{qoptions}
\opt{A}{Amount of heat generated in the circuit as the capacitor gets charged is $\dfrac12CV^2$}
\opt{B}{Magnitude of force that one capacitor plate exerts on the other is $\dfrac{\varepsilon_0AV^2}{d^2}$}
\opt{C}{If the distance between the plates is slowly reduced to $\dfrac d2$, the work done by external agent in the process is given by $-\dfrac12CV^2$}
\opt{D}{If the distance between the plates is reduced slowly from $d$ to $\dfrac d2$, energy dissipated in the circuit is zero.}
\end{qoptions}
\end{bankq}
\begin{bankq}[id=TOAS11-P2-PHY-S3-Q03, source={TOAS - Test 11 (27 Jun 2022)}, paper=2, qno=3, type=numerical, subtopic={RC circuits}, answer={1.00}, answer_src=key, difficulty=4, diff_parts={K1 S2 I1 C0 F1}, diff_note={concepts: RC charging with a resistor network | steps: Thevenin equivalent seen by the capacitor; final charge; time constant; alpha | I: Thevenin equivalent}]
Write the variation in charge ($q$) on the capacitor plates as a function of time after the switch (S) is closed at $t=0$. If $q$ is given by $q=\dfrac{CV}{2}\left[1-e^{-\frac{2t}{\alpha RC}}\right]$, then find the value of $\alpha$.
\bfig{106.3}{TOAS11_p2_phy_s3_q03_1.png}
\end{bankq}
\begin{bankq}[id=TOAS12-P1-PHY-S1-Q04, source={TOAS - Test 12 (10 Jul 2022)}, paper=1, qno=4, type=multi, subtopic={Capacitors}, answer={B,C}, answer_src=key, difficulty=3, diff_parts={K1 S2 I0 C0 F1}, diff_note={concepts: capacitor with a constant current source | steps: dV/dt = I/C; voltage after 10 ms; sign from the given direction}]
Diagram shows a capacitor connected with a constant current source. Choose the correct option(s) :
\bfig{66.6}{TOAS12_p1_phy_s1_q04_1.png}
\begin{qoptions}
\opt{A}{the rate of change of voltage across the capacitor in the circuit of figure from top to bottom 166.7 V/s.}
\opt{B}{the rate of change of voltage across the capacitor in the circuit of figure from top to bottom $-$166.7 V/s.}
\opt{C}{the capacitor's voltage 10 milliseconds after power is switched on is 1.667V.}
\opt{D}{the capacitor's voltage 10 milliseconds after power is switched on is 2.667V.}
\end{qoptions}
\end{bankq}
\begin{bankq}[id=TOAS12-P1-PHY-S1-Q06, source={TOAS - Test 12 (10 Jul 2022)}, paper=1, qno=6, type=multi, subtopic={Energy}, answer={A,C}, answer_src=key, difficulty=5, diff_parts={K2 S2 I1 C0 F1}, diff_note={concepts: capacitor with a conducting disc, electrostatic force | steps: capacitance chi of the disc with the plate; force from the energy; threshold voltage from the weight | I: force from energy at constant voltage}]
A capacitor consists of two circular parallel plates both with radius $R$ separated by distance $d$, where $d\ll R$, as shown in figure (a). The top plate is connected to a constant voltage source at a potential $V$ while the bottom plate is grounded. Then a thin and small disk of mass $m$ with radius $r$ ($\ll R,d$) and thickness $t$ ($\ll r$) is placed on the center of the bottom plate, as shown in figure (b).\par
Let us assume that the space between the plates is in vacuum with the dielectric constant $\epsilon_0$; the plates and the disk are made of perfect conductors; and all the electrostatic edge effects may be neglected. The inductance of the whole circuit and the relativistic effects can be safely disregarded. The image charge effect can also be neglected. Choose the correct option(s)
\bfig{210.6}{TOAS12_p1_phy_s1_q06_1.png}
\begin{qoptions}
\opt{A}{When the disk is placed on the bottom plate, a charge $q$ on the disk of figure (b) is related to the voltage $V$ by $q=\chi V$, then $\chi$ in terms of $r$, $d$ and $\varepsilon_0$ is given by $-\dfrac{\epsilon_0\pi r^2}{d}$.}
\opt{B}{When the disk is placed on the bottom plate, a charge $q$ on the disk of figure (b) is related to the voltage $V$ by $q=\chi V$, then $\chi$ in terms of $r$, $d$ and $\varepsilon_0$ is given by $\dfrac{\epsilon_0\pi r^2}{d}$.}
\opt{C}{The parallel plates lie perpendicular to a uniform gravitational field $g$. To lift up the disk at rest initially, we need to increase the applied voltage beyond a threshould voltage $v_{th}$. The value of $v_{th}$ is given by $\sqrt{\dfrac{2mgd}{\chi}}$.}
\opt{D}{The parallel plates lie perpendicular to a uniform gravitational field $g$. To lift up the disk at rest initially, we need to increase the applied voltage beyond a threshould voltage $v_{th}$. The value of $v_{th}$ is given by $\sqrt{\dfrac{mgd}{\chi}}$.}
\end{qoptions}
\end{bankq}
\begin{bankq}[id=TOAS12-P1-PHY-S2-Q03, source={TOAS - Test 12 (10 Jul 2022)}, paper=1, qno=3, type=numerical, subtopic={Dielectrics}, answer={1.50}, answer_src=key, difficulty=3, diff_parts={K1 S2 I0 C0 F1}, diff_note={concepts: capacitor with two dielectrics | steps: series capacitors; the field in each slab; ratio from the potential difference}]
A parallel plate capacitor consisting two dielectric materials is shown in the figure. The middle dielectric slab is place symmetrically with respect to the plates. If the potential difference between one of the plates and the nearest surface of dielectric interface is 2Volts, then the ratio $\epsilon_1:\epsilon_2$ is
\bfig{103.8}{TOAS12_p1_phy_s2_q03_1.png}
\end{bankq}
\begin{bankq}[id=TOAS12-P2-PHY-S1-Q06, source={TOAS - Test 12 (10 Jul 2022)}, paper=2, qno=6, type=single, subtopic={RC circuits}, answer={D}, answer_src=key, difficulty=2, diff_parts={K1 S2 I0 C0 F0}, diff_note={concepts: RC circuit with an initially charged capacitor | steps: initial current from loop equation; final charges on C and C0; effect of q0}]
In the circuit shown, the switch S is open initially. The capacitor $C$ is uncharged but $C_0$ has a charge $q_0$ at $t=0$
\bfig{109.0}{TOAS12_p2_phy_s1_q06_1.png}
\begin{qoptions}
\opt{A}{The current immediately after switch S is closed is $I_C=\dfrac{\left(E+\frac{q_0}{C_0}\right)}{R}$.}
\opt{B}{The current immediately after switch S is closed is $I_c=\dfrac{E}{R}$}
\opt{C}{Final charge on $C$ after closing switch is $\dfrac{EC_0C}{C_0+C}$}
\opt{D}{Final charge on $C$ after closing switch is $\dfrac{EC_0C}{C_0+C}\left(1-\dfrac{q_0}{C_0E}\right)$}
\end{qoptions}
\end{bankq}
\begin{bankq}[id=TOAS12-P2-PHY-S3-Q01, source={TOAS - Test 12 (10 Jul 2022)}, paper=2, qno=1, type=numerical, subtopic={RC circuits}, answer={81.00}, answer_src=key, difficulty=3, diff_parts={K1 S2 I0 C0 F1}, diff_note={concepts: charging of capacitors through a resistor | steps: charge on A at half the final value; time; heat as energy from the source minus stored energy}]
In the circuit shown in the figure, find the heat generated (in $\mu$J) through the resistor upto the instant when charge on capacitor A is half of the maximum values that the capacitors would store finally. (Initially both the capacitor are uncharged) (Given : $C=8\ \mu$F, $\varepsilon=3$ volt).
\bfig{65.9}{TOAS12_p2_phy_s3_q01_1.png}
\end{bankq}
\begin{bankq}[id=TOAS12-P2-PHY-S3-Q02, source={TOAS - Test 12 (10 Jul 2022)}, paper=2, qno=2, type=numerical, subtopic={Capacitors}, answer={9.00}, answer_src=key, difficulty=4, diff_parts={K1 S2 I1 C0 F1}, diff_note={concepts: capacitor with a moving middle plate | steps: series capacitors C1 and C2 with the original gaps; new gaps after the shift; potential at the middle node | I: potential divider of capacitors}]
In a gravity free space, three identical metal plates have been placed parallel to each other as shown. In the initial state shown in the diagram, the potential difference between the points A and B is $V_1$. Now, the middle plate is shifted down a distance $L/2$, keeping the battery connected. The other plates are kept fixed. The potential difference between the points A and B now becomes $V_2$. The value $\dfrac{6V_1}{V_2}$ is equal to
\bfig{112.0}{TOAS12_p2_phy_s3_q02_1.png}
\end{bankq}
\begin{bankpassage}[id=TOAS13-P1-PHY-S2-G9-10, source={TOAS - Test 13 (24 Jul 2022)}, questions={9-10}]
Consider the circuit shown in figure. Battery is ideal and capacitor was initially \underline{uncharged}. If at $t=0$ switch is closed, then answer the following question :
\bfig{109.4}{TOAS13_p1_phy_s2_para9_10_1.png}
\end{bankpassage}
\begin{bankq}[id=TOAS13-P1-PHY-S2-Q09, source={TOAS - Test 13 (24 Jul 2022)}, paper=1, qno=9, type=multi, subtopic={RC circuits}, answer={B}, answer_src=key, difficulty=3, diff_parts={K1 S2 I0 C0 F1}, diff_note={concepts: RC circuit at t = 0 | steps: uncharged capacitor acts as a short; reduce the resistor network; current in the marked branch with its sign}, group=TOAS13-P1-PHY-S2-G9-10]
Current in the direction B shown in diagram at $t=0$ (Just after the switch is closed) will be :-
\begin{qoptions}
\opt{A}{$\dfrac38$ A}
\opt{B}{$\left(\dfrac{-3}{8}\right)$ A}
\opt{C}{$\left(\dfrac45\right)$ A}
\opt{D}{$\left(\dfrac{-4}{5}\right)$ A}
\end{qoptions}
\end{bankq}
\begin{bankq}[id=TOAS13-P1-PHY-S2-Q10, source={TOAS - Test 13 (24 Jul 2022)}, paper=1, qno=10, type=multi, subtopic={RC circuits}, answer={C}, answer_src=key, difficulty=5, diff_parts={K2 S2 I0 C1 F1}, diff_note={concepts: Thevenin resistance, RC transient | steps: resistance seen by the capacitor; time constant tau = R C; initial current; I = I0 exp(-t/tau)}, group=TOAS13-P1-PHY-S2-G9-10]
Current $I$ as a function of $t$ (in $\mu$s) will be :-
\begin{qoptions}
\opt{A}{$I=\dfrac94e^{-3t/8}$}
\opt{B}{$I=\dfrac32e^{-3t/10}$}
\opt{C}{$I=\dfrac94e^{-3t/16}$}
\opt{D}{$I=\dfrac32e^{-3t/5}$}
\end{qoptions}
\end{bankq}
\begin{bankq}[id=TOAS13-P2-PHY-S1-Q08, source={TOAS - Test 13 (24 Jul 2022)}, paper=2, qno=8, type=single, subtopic={Combinations}, answer={B}, answer_src=key, difficulty=2, diff_parts={K1 S2 I0 C0 F0}, diff_note={concepts: series and parallel capacitors | steps: equivalent capacitance; total charge; split between branches}]
In the circuit shown in figure.
\bfig{121.7}{TOAS13_p2_phy_s1_q08_1.png}
\begin{qoptions}
\opt{A}{$q_2=280\mu$C}
\opt{B}{$q_2=360\mu$C}
\opt{C}{$q_3=360\mu$C}
\opt{D}{It is not possible to find $q_2$ and $q_3$ unless $C_1$ and $C_2$ are known.}
\end{qoptions}
\end{bankq}
\begin{bankq}[id=TOAS13-P2-PHY-S1-Q15, source={TOAS - Test 13 (24 Jul 2022)}, paper=2, qno=15, type=single, subtopic={Capacitors}, answer={C}, answer_src=key, difficulty=1, diff_parts={K1 S1 I0 C0 F0}, diff_note={concepts: current through a capacitor | steps: I = C dV/dt from the graph slopes gives C = 2 F}]
Voltage versus time and current versus time graph of a circuit element is shown in figure. The type of circuit element is :
\bfig{152.8}{TOAS13_p2_phy_s1_q15_1.png}
\begin{qoptions}
\opt{A}{resistance of 2$\Omega$}
\opt{B}{capacitance of 1F}
\opt{C}{capacitance of 2F}
\opt{D}{none}
\end{qoptions}
\end{bankq}
\begin{bankq}[id=TOAS13-P2-PHY-S2-Q02, source={TOAS - Test 13 (24 Jul 2022)}, paper=2, qno=2, type=numerical, subtopic={Charge sharing}, answer={120.00}, answer_src=key, difficulty=3, diff_parts={K1 S2 I0 C0 F1}, diff_note={concepts: capacitor charge flow after closing a switch | steps: charges before closing; charges after closing; difference at the node}]
What amount of charge (in $\mu$C) will flow through switch S from A to B, when switch S is closed.
\bfig{132.5}{TOAS13_p2_phy_s2_q02_1.png}
\end{bankq}
\begin{bankq}[id=TOAS14-P1-PHY-S1-Q08, source={TOAS - Test 14 (07 Aug 2022)}, paper=1, qno=8, type=multi, subtopic={Dielectrics}, answer={A,B,D}, answer_src=key, difficulty=4, diff_parts={K1 S2 I0 C1 F1}, diff_note={concepts: capacitors with dielectrics between four plates | steps: identify the series and parallel pieces; equivalent capacitance; plate charges from the applied voltage; energy}]
Four plates of equal area $A$ and plate seperation $d$ are arranged as shown and filled with different dielectric having dielectric constant $K$. $\left(\dfrac{\epsilon_0A}{d}=1\mu\text{F}\right)$ Select correct alternative.
\bfig{134.7}{TOAS14_p1_phy_s1_q08_1.png}
\begin{qoptions}
\opt{A}{Effective capacitance of system is $\dfrac43\mu$F .}
\opt{B}{Charge on $2^{\text{nd}}$ plate is $-10~\mu$C.}
\opt{C}{Charge on $3^{\text{rd}}$ plate is $+\dfrac{40}{3}\mu$C.}
\opt{D}{Energy stored in system is $\dfrac{200}{3}\mu$J.}
\end{qoptions}
\end{bankq}
\begin{bankq}[id=TOAS15-P1-PHY-S2-Q07, source={TOAS - Test 15 (22 Aug 2022)}, paper=1, qno=7, type=multi, subtopic={Combinations}, answer={B,C}, answer_src=key, difficulty=5, diff_parts={K2 S2 I0 C1 F1}, diff_note={concepts: capacitance of concentric shells, dielectric | steps: capacitance of each pair of shells; the outer shell also has capacitance to infinity; parallel combination; charge on the middle shell from V = Q/C}]
Three concentric conducting spherical shells have radius $R$, $2R$ \& $3R$ are shown. Middle shell is earthed and inner most shell is connected to outer most shell by conducting wire. Space between $2^{\text{nd}}$ and $3^{\text{rd}}$ shell is filled by dielectric of dielectric constant $K=2$ and $Q=17\mu$C charge is given to $3^{\text{rd}}$ shell then :-
\bfig{84.6}{TOAS15_p1_phy_s2_q07_1.png}
\begin{qoptions}
\opt{A}{Capacitance of system is $44\pi\epsilon_0R$}
\opt{B}{Capacitance of system is $68\pi\epsilon_0R$}
\opt{C}{Charge on $2^{\text{nd}}$ shell is $-14\mu$C.}
\opt{D}{Charge on $2^{\text{nd}}$ shell is $-17\mu$C.}
\end{qoptions}
\end{bankq}
\begin{bankq}[id=TOAS15-P2-PHY-S2-Q03, source={TOAS - Test 15 (22 Aug 2022)}, paper=2, qno=3, type=numerical, subtopic={Dielectrics}, answer={8.46}, answer_src=key, difficulty=2, diff_parts={K1 S1 I0 C0 F1}, diff_note={concepts: breakdown field and capacitance | steps: minimum plate separation from the voltage and field; K from C = K eps0 A/d}]
Voltage rating of a parallel plate capacitor is 500 V. Its dielectric can withstand a maximum electric field of $10^6$ V/m. The plate area is $10^{-4}$m$^2$. What is the dielectric constant if the capacitance is 15pF ? (given $\epsilon_0=8.86\times10^{-12}$ C$^2$ m$^2$)
\end{bankq}
\begin{bankq}[id=TOAS15-P2-PHY-S3-Q04, source={TOAS - Test 15 (22 Aug 2022)}, paper=2, qno=4, type=numerical, subtopic={Energy}, answer={3}, answer_src=key, difficulty=5, diff_parts={K2 S2 I1 C0 F1}, diff_note={concepts: energy of capacitors with, without the battery | steps: energy before; capacitor on the battery keeps V; isolated capacitor keeps its charge; ratio}]
The figure shows two identical parallel plate capacitors connected to a battery with switch S closed. The switch is now opened and the free space between the plates of the capacitors is filled with a dielectric of dielectric constant (or relative permittivity) 3. The ratio of the total electrostatic energy stored in both capacitors before and after the introduction of the dielectric is $\dfrac x5$. Then find the value of $x$
\bfig{99.9}{TOAS15_p2_phy_s3_q04_1.png}
\end{bankq}
\begin{bankq}[id=TOAS18-P1-PHY-S1-Q05, source={TOAS - Test 18 (30 Oct 2022)}, paper=1, qno=5, type=multi, subtopic={Dielectrics}, answer={A,B}, answer_src=solved, working={Dielectric occupies angle $\pi-\phi$: $C\propto\epsilon-\tfrac\phi\pi(\epsilon-\epsilon_0)$; at fixed $V$ torque $=\tfrac12V^2dC/d\phi=V^2R^2(\epsilon-\epsilon_0)/4d$ (constant, so no SHM); the expression in C is not the torque at fixed charge}, difficulty=5, diff_parts={K2 S2 I1 C0 F1}, diff_note={concepts: capacitance with a rotating dielectric, energy method | steps: capacitance as a function of the angle; torque at fixed voltage; behaviour of the torque; fixed charge case}]
A capacitor consists of two stationary plates shaped as a semi-circle of radius $R$ and a movable plate made of dielectric with permittivity $\in$ and capable of rotating about an axis $O$ between the stationary plates. The thickness of movable plate is $d$, which is nearly same as separation between stationary plates. Select correct statement(s). $\phi$ is angle of sector outside of capacitor.
\bfig{87.2}{TOAS18_p1_phy_s1_q05_1.png}
\begin{qoptions}
\opt{A}{Capacitance of system is proportional to $\left\{\in-\dfrac\phi\pi(\in-\in_0)\right\}$}
\opt{B}{If potential difference $V$ is applied to capacitor electric torque about axis $O$ is $\dfrac{V^2R^2}{4d}(\in-\in_0)$ acting on movable plate.}
\opt{C}{If this capacitor is charged by charge $Q$ and then source is removed then electric torque about axis $O$ is $\dfrac{Q^2}{4(\in-\in_0)\,d}$.}
\opt{D}{If potential difference $V$ is applied to capacitor and system is released than movable plate will perform simple harmonic motion.}
\end{qoptions}
\end{bankq}
\begin{bankq}[id=TOAS18-P2-PHY-S3-Q06, source={TOAS - Test 18 (30 Oct 2022)}, paper=2, qno=6, type=numerical, subtopic={Combinations}, answer={1}, answer_src=key, difficulty=4, diff_parts={K1 S2 I0 C1 F1}, diff_note={concepts: capacitor energy for three plates | steps: capacitances; equivalent; energy}]
3 plates $A$, $B$ \& $C$ each of area 1 m$^2$ are separated by 0.885 mm. Energy stored in system (in $\mu$J) is
\bfig{105.3}{TOAS18_p2_phy_s3_q06_1.png}
\end{bankq}

\section*{TAAS papers}
\begin{bankq}[id=TAAS01-P1-PHY-S1-Q02, source={TAAS - Test 01 (18 Jun 2023)}, paper=1, qno=2, type=single, subtopic={Dielectrics}, answer={D}, answer_src=key, difficulty=3, diff_parts={K1 S2 I1 C0 F0}, diff_note={concepts: bound charge at a dielectric boundary | steps: displacement is continuous; field in each layer; polarisation difference gives the bound charge density}]
The space between the plates of a parallel plate capacitor is filled completely consecutively with two dielectric layers 1 and 2 having the thickness $d_1$ and $d_2$ and dielectric constant $k_1$ and $k_2$ respectively. The area of each plate is equal to $A$. The density $\sigma'$ of the bound charges on the boundary plane between the two dielectrics if the voltage across the capacitor equals $V$ and electric field is directed from layer 1 to layer 2 is :
\begin{qoptions}
\opt{A}{$\left(\dfrac{k_1-k_2}{k_1+k_2}\right)\dfrac{\varepsilon_0V}{(d_1+d_2)}$}
\opt{B}{$\left(\dfrac{k_1-k_2}{k_1d_1+k_2d_2}\right)\varepsilon_0V$}
\opt{C}{$\left(\dfrac{k_2d_2-k_1d_1}{k_1+k_2}\right)\dfrac{\varepsilon_0V}{d_1d_2}$}
\opt{D}{$\left(\dfrac{k_1-k_2}{k_1d_2+k_2d_1}\right)\varepsilon_0V$}
\end{qoptions}
\end{bankq}
\begin{bankq}[id=TAAS01-P1-PHY-S2-Q09, source={TAAS - Test 01 (18 Jun 2023)}, paper=1, qno=9, type=multi, subtopic={Dielectrics}, answer={A,D}, answer_src=key, difficulty=4, diff_parts={K1 S2 I1 C0 F1}, diff_note={concepts: potential and field in a half-filled capacitor | steps: field in each half; potential variation; plot shapes}]
In the adjacent figure, half of the space between the plates of a parallel plate capacitor is filled with dielectric material of constant $K$. Then which of the plots are possible if $X$ is measured from the plate $A$ :- ($V=$ potential, $E=$ electric field)
\bfig{99.7}{TAAS01_p1_phy_s2_q09_1.png}
\begin{qoptions}
\opt{A}{\bfigi{65.9}{TAAS01_p1_phy_s2_q09_A.png}}
\opt{B}{\bfigi{66.1}{TAAS01_p1_phy_s2_q09_B.png}}
\opt{C}{\bfigi{65.9}{TAAS01_p1_phy_s2_q09_C.png}}
\opt{D}{\bfigi{55.6}{TAAS01_p1_phy_s2_q09_D.png}}
\end{qoptions}
\end{bankq}
\begin{bankq}[id=TAAS01-P1-PHY-S2-Q10, source={TAAS - Test 01 (18 Jun 2023)}, paper=1, qno=10, type=multi, subtopic={RC circuits}, answer={A,C}, answer_src=key, difficulty=3, diff_parts={K1 S2 I0 C0 F1}, diff_note={concepts: voltages in an RC circuit | steps: sum of resistor and capacitor voltages is constant; identify the curves; relate the plotted voltages}]
In the given $R-C$ circuit switch $S$ is closed at $t=0$ and the variation of potential drop across the resistor ($V_R$) and across capacitor ($V_C$) is represented by given plots, then:
\bfig{136.2}{TAAS01_p1_phy_s2_q10_1.png}
\begin{qoptions}
\opt{A}{$2V_0=2V_1+V_2$}
\opt{B}{$2V_0=V_1+V_2$}
\opt{C}{Curve I is for $V_C$ and II is for $V_R$.}
\opt{D}{Curve I is for $V_R$ and II is for $V_C$.}
\end{qoptions}
\end{bankq}
\begin{bankq}[id=TAAS02-P1-PHY-S1-Q02, source={TAAS - Test 02 (30 Jul 2023)}, paper=1, qno=2, type=single, subtopic={Combinations}, answer={A}, answer_src=key, difficulty=4, diff_parts={K1 S2 I1 C1 F0}, diff_note={concepts: equivalent capacitance of a bridge network | steps: identify balanced or unbalanced bridge; reduce series and parallel parts; combine; value}]
In the circuit diagram shown all the capacitors are in $\mu$F. The equivalent capacitance (in $\mu$F) between points $A$ and $B$ is :-
\bfig{108.3}{TAAS02_p1_phy_s1_q02_1.png}
\begin{qoptions}
\opt{A}{$\dfrac{14}5$}
\opt{B}{$\dfrac75$}
\opt{C}{$\dfrac37$}
\opt{D}{None of these}
\end{qoptions}
\end{bankq}
\begin{bankq}[id=TAAS02-P1-PHY-S1-Q03, source={TAAS - Test 02 (30 Jul 2023)}, paper=1, qno=3, type=single, subtopic={Combinations}, answer={C}, answer_src=key, difficulty=4, diff_parts={K1 S2 I1 C1 F0}, diff_note={concepts: steady-state charges with resistors | steps: steady state currents; voltages across the capacitors; charges; ratio}]
Ratio of charges on 5$\mu$F and 10$\mu$F capacitor when circuit is in steady state :-
\bfig{107.5}{TAAS02_p1_phy_s1_q03_1.png}
\begin{qoptions}
\opt{A}{38 : 13}
\opt{B}{19 : 26}
\opt{C}{13 : 38}
\opt{D}{26 : 19}
\end{qoptions}
\end{bankq}
\begin{bankq}[id=TAAS02-P1-PHY-S4-Q02, source={TAAS - Test 02 (30 Jul 2023)}, paper=1, qno=2, type=numerical, subtopic={Combinations}, answer={4}, answer_src=key, difficulty=5, diff_parts={K1 S2 I1 C1 F1}, diff_note={concepts: series capacitors with a parallel addition | steps: voltage division gives the capacitance ratio; new voltages after adding 2 microfarad; solve}]
Two capacitors A and B are connected in series across a 100V supply and it is observed that the potential difference across them are 60V and 40V. A capacitor of 2$\mu$F capacitance is now connected in parallel with A and the potential difference across B raises to 90 V. If capacitance of A is 0.04 x$\mu$F. Find $x$.
\end{bankq}
\begin{bankq}[id=TAAS02-P2-PHY-S2-Q07, source={TAAS - Test 02 (30 Jul 2023)}, paper=2, qno=7, type=multi, subtopic={Energy}, answer={A,C,D}, answer_src=key, difficulty=5, diff_parts={K1 S2 I1 C1 F1}, diff_note={concepts: charging three capacitors | steps: equivalent capacitance; charge flow; energy supplied and stored; heat}]
Three capacitors each having capacitance $C=2\,\mu$F are connected with a battery of emf 30 V as shown in figure. Select the correct statement(s), when the switch $S$ is closed.
\bfig{58.5}{TAAS02_p2_phy_s2_q07_1.png}
\begin{qoptions}
\opt{A}{The amount of charge flown through the battery is 20 $\mu$C.}
\opt{B}{The heat generated in the circuit is 0.6 mJ.}
\opt{C}{The energy supplied by the battery is 0.6 mJ.}
\opt{D}{The amount of charge flown through the switch $S$ is 60 $\mu$C.}
\end{qoptions}
\end{bankq}

\section*{TEAS papers}
\begin{bankq}[id=TEAS01-P1-PHY-S1-Q04, source={TEAS - Test 01 (16 Apr 2023)}, paper=1, qno=4, type=multi, subtopic={Combinations}, answer={A,B,C}, answer_src=key, difficulty=5, diff_parts={K1 S2 I1 C1 F1}, diff_note={concepts: capacitors with batteries | steps: charge on the series chain; loop equation; voltage on each capacitor; check}]
In the diagram we have three capacitors $C_1=20\ \mu$F, $C_2=30\ \mu$F and $C_3=60\ \mu$F, two sources with emf $\varepsilon_1=1$ V and $\varepsilon_2=2$ V. The potential difference $V_A-V_B=3$ V. Find the voltage on each capacitor.
\bfig{161.3}{TEAS01_p1_phy_s1_q04_1.png}
\begin{qoptions}
\opt{A}{$V_1=1$ volt}
\opt{B}{$V_2=2/3$ volt}
\opt{C}{$V_3=1/3$ volt}
\opt{D}{$V_1=2$ volt}
\end{qoptions}
\end{bankq}

\section*{TMAS papers}
\begin{bankq}[id=TMAS02-P1-PHY-S1-Q01, source={TMAS - Test 02 (06 Aug 2023)}, paper=1, qno=1, type=single, subtopic={Dielectrics}, answer={D}, answer_src=key, difficulty=3, diff_parts={K1 S1 I1 C1 F0}, diff_note={concepts: capacitor with dielectric sections | steps: identify series and parallel parts; combine}]
An air capacitor having capacity $C_0$ is filled with dielectric as shown in figure. then the new capacitance of system will be
\bfig{129.4}{TMAS02_p1_phy_s1_q01_1.png}
\begin{qoptions}
\opt{A}{$C_0$}
\opt{B}{$2C_0$}
\opt{C}{$3C_0$}
\opt{D}{$4C_0$}
\end{qoptions}
\end{bankq}
\begin{bankq}[id=TMAS02-P1-PHY-S1-Q02, source={TMAS - Test 02 (06 Aug 2023)}, paper=1, qno=2, type=single, subtopic={Combinations}, answer={C}, answer_src=key, difficulty=5, diff_parts={K2 S2 I1 C1 F0}, diff_note={concepts: series capacitors, graph reading | steps: voltage divider; read the graph; solve the two unknowns}]
Capacitor 3 in circuit shown below is a variable capacitor. Graph shows the electric potential across capacitor 1 versus $C_3$. The electric potential $V$ across the battery and capacitance $C_1$ is :-
\bfig{190.6}{TMAS02_p1_phy_s1_q02_1.png}
\begin{qoptions}
\opt{A}{$2V,\ 8\ \mu F$}
\opt{B}{$10V,\ 6\ \mu F$}
\opt{C}{$10V,\ 8\ \mu F$}
\opt{D}{$10V,\ 12\ \mu F$}
\end{qoptions}
\end{bankq}
\begin{bankq}[id=TMAS02-P1-PHY-S2-Q05, source={TMAS - Test 02 (06 Aug 2023)}, paper=1, qno=5, type=multi, subtopic={RC circuits}, answer={B,C,D}, answer_src=key, difficulty=5, diff_parts={K1 S2 I1 C1 F1}, diff_note={concepts: steady state of an RC circuit | steps: capacitor open; current path; voltage and energy}]
In the circuit shown in figure the switch $S$ is closed at $t=0$. A long time after closing the switch
\bfig{116.5}{TMAS02_p1_phy_s2_q05_1.png}
\begin{qoptions}
\opt{A}{Voltage drop across the capacitor is $E$}
\opt{B}{Current through the battery is $\dfrac{E}{R_1+R_2+R_3}$}
\opt{C}{Energy stored in the capacitor is {\footnotesize$\dfrac12C\!\left(\dfrac{(R_2+R_3)E}{R_1+R_2+R_3}\right)^2$}}
\opt{D}{Current through the resistance $R_4$ becomes zero}
\end{qoptions}
\end{bankq}
\begin{bankq}[id=TMAS02-P1-PHY-S2-Q07, source={TMAS - Test 02 (06 Aug 2023)}, paper=1, qno=7, type=multi, subtopic={RC circuits}, answer={A,B,D}, answer_src=key, difficulty=5, diff_parts={K1 S2 I1 C1 F1}, diff_note={concepts: capacitor discharging into another | steps: time constant; current; charge; heat loss}]
In the circuit shown in figure $C_1=C_2=C$ and capacitor $C_1$ is having initial charge $Q_0$. The switch is closed at $t=0$. Then choose the correct option(s):
\bfig{102.1}{TMAS02_p1_phy_s2_q07_1.png}
\begin{qoptions}
\opt{A}{The current at time $t$ in the circuit will be $\left(\dfrac{Q_0}{RC}\right)e^{-2t/RC}$}
\opt{B}{Charge on capacitor $C_1$ at time $t$ will be $\dfrac{Q_0}{2}\left(1+e^{-2t/RC}\right)$}
\opt{C}{Total heat loss in the circuit till the current becomes zero is $\dfrac{Q_0^2}{2C}$}
\opt{D}{Total heat loss in the circuit till the current becomes zero is $\dfrac{Q_0^2}{4C}$}
\end{qoptions}
\end{bankq}
\begin{bankq}[id=TMAS02-P2-PHY-S1-Q03, source={TMAS - Test 02 (06 Aug 2023)}, paper=2, qno=3, type=multi, subtopic={RC circuits}, answer={A,C,D}, answer_src=solved, working={Steady state is a single series loop of six $10\ \Omega$ resistors with net emf $80-20=60$ V, $I=1$ A. Then $V_A-V_B=40$ V, $|Q_1|=30$ C, $|Q_2|=50$ C}, difficulty=6, diff_parts={K1 S2 I2 C1 F1}, diff_note={concepts: DC circuit with capacitors | steps: steady state loop; net emf; potentials; charges}]
For the given circuit in steady state, choose the correct statements :- (given all capacitors have capacitance = $1$ in SI \& all resistors have resistance = $10$ in SI unit)
\bfig{108.3}{TMAS02_p2_phy_s1_q03_1.png}
\begin{qoptions}
\opt{A}{Potential difference between point $A$ \& $B$ $(V_A-V_B)$ is $40$ V.}
\opt{B}{Potential difference between point $A$ \& $B$ $(V_A-V_B)$ is $50$ V.}
\opt{C}{Magnitude of charge $Q_1$ is $30$ C.}
\opt{D}{Magnitude of charge $Q_2$ is $50$ C.}
\end{qoptions}
\end{bankq}
\begin{bankq}[id=TMAS02-P2-PHY-S1-Q04, source={TMAS - Test 02 (06 Aug 2023)}, paper=2, qno=4, type=multi, subtopic={Dielectrics}, answer={B,D}, answer_src=solved, working={$C=10\varepsilon_0A/d=8.85$ nF gives $d=1$ cm; breakdown limit is $10$ kV; a 0.001 rad tilt over 1 m changes the gap by 1 mm out of 10 mm, far worse than a 0.1 percent error}, difficulty=7, diff_parts={K2 S2 I2 C1 F1}, diff_note={concepts: capacitor design, tilted plates | steps: capacitance formula; plate separation; breakdown voltage; sensitivity}]
As preparation for the olympiads, Ivan wanted to make a graphite-filled, parallel plate, $8.85$ nF capacitor with square plates of side $1$ m. ``Such big metal plates will be expensive,'' said his father. ``Can you manage with cardboard? Cover it with aluminium foil to have metal.'' His sister who studies physics in college, said, ``To get the correct capacitance'', make sure you keep the plates parallel. Having parallel plates with a fractional error of $0.001$ in plate separation is better than having the correct plate separation at one end and an angle of $0.001$ radians between the plates.'' The dielectric constant of graphite is $10$ and its dielectric strength is $1$ kV/mm. Which of the following statements are correct?
\begin{qoptions}
\opt{A}{The aluminium foil on cardboard will not work.}
\opt{B}{The plate separation required is $1$ cm.}
\opt{C}{The capacitor can withstand a maximum potential difference of $1000$ V.}
\opt{D}{Ivan's sister is right.}
\end{qoptions}
\end{bankq}
\begin{bankpassage}[id=TMAS02-P2-PHY-S2-G9-10, source={TMAS - Test 02 (06 Aug 2023)}, questions={9-10}]
At $t=0$, switch is closed. Capacitors are initially uncharged. Current through battery at $t=0^+$ is $(10\alpha)$ Ampere while that at $t\to\infty$ is $(10\beta)$ Ampere.
\bfig{158.5}{TMAS02_p2_phy_s2_para9_10_1.png}
\end{bankpassage}
\begin{bankq}[id=TMAS02-P2-PHY-S2-Q09, source={TMAS - Test 02 (06 Aug 2023)}, paper=2, qno=9, type=single, subtopic={RC circuits}, answer={B}, answer_src=solved, working={At $t=0^+$ capacitors are short circuits, equivalent resistance $4\ \Omega$ including $2\ \Omega$, so $I=15$ A and $\alpha=1.5$}, difficulty=5, diff_parts={K1 S2 I2 C1 F0}, diff_note={concepts: circuit at the moment of switching | steps: capacitors as shorts; equivalent resistance; current}, group=TMAS02-P2-PHY-S2-G9-10]
The value of $\alpha$ is \underline{\hspace{1.5cm}}
\begin{qoptions}
\opt{A}{$2$}
\opt{B}{$1.5$}
\opt{C}{$4.5$}
\opt{D}{$1$}
\end{qoptions}
\end{bankq}
\begin{bankq}[id=TMAS02-P2-PHY-S2-Q10, source={TMAS - Test 02 (06 Aug 2023)}, paper=2, qno=10, type=single, subtopic={RC circuits}, answer={A}, answer_src=solved, working={At steady state resistance is $10+8+2=20\ \Omega$, $I=3$ A, so $\beta=0.3$}, difficulty=5, diff_parts={K1 S2 I2 C1 F0}, diff_note={concepts: steady state circuit | steps: capacitors open; remaining resistors; current}, group=TMAS02-P2-PHY-S2-G9-10]
Value of $\beta$ is \underline{\hspace{1.2cm}}
\begin{qoptions}
\opt{A}{$0.3$}
\opt{B}{$3$}
\opt{C}{$1$}
\opt{D}{$2$}
\end{qoptions}
\end{bankq}

\section*{TNAS papers}
\begin{bankq}[id=TNAS02-P1-PHY-S2-Q10, source={TNAS - Test 02 (07 May 2023)}, paper=1, qno=10, type=multi, subtopic={Dielectrics}, answer={A,B,C}, answer_src=key, difficulty=7, diff_parts={K1 S3 I2 C1 F1}, diff_note={concepts: capacitor with dielectric slab and block | steps: capacitance at half insertion; charge; energy conservation; kinetic energy; check}]
A capacitor with air as dielectric has capacity $C$ as shown in figure is completely filled with dielectric slab of mass $M$ and is connected to a block of mass $M$. The system is released form rest at time $t=0$ when the capacitor is completely filled with dielectric. Then: (Dielectric constant of slab $=K$)
\bfig{121.7}{TNAS02_p1_phy_s2_q10_1.png}
\begin{qoptions}
\opt{A}{Charge on the plate of the capacitor when it is half filled with air is $\dfrac{(K+1)CV}{2}$}
\opt{B}{The kinetic energy of the block when the capacitor is half filled with dielectric is $\dfrac{mg\ell}{4}-\dfrac{CV^2(K-1)}{8}$}
\opt{C}{The potential difference between the plates when the capacitor is half filled with air is $V$.}
\opt{D}{The energy stored in the capacitor when it is half filled with dielectric is $K/2\,CV^2$.}
\end{qoptions}
\end{bankq}
\begin{bankpassage}[id=TNAS02-P1-PHY-S3-G3-4, source={TNAS - Test 02 (07 May 2023)}, questions={3-4}]
In the given circuit diagram, both capacitors are initially uncharged. The capacitance $C_1=2\,\mu$F and $C_2=4\,\mu$F, emf of battery A and B are 2V and 4V respectively.
\bfig{108.5}{TNAS02_p1_phy_s3_para3_4_1.png}
\end{bankpassage}
\begin{bankq}[id=TNAS02-P1-PHY-S3-Q03, source={TNAS - Test 02 (07 May 2023)}, paper=1, qno=3, type=numerical, subtopic={Energy}, answer={8}, answer_src=solved, working={S1 closed: C1 across battery A only, $Q=2\times2=4\ \mu$C; work by battery $=QV=8\ \mu$J}, difficulty=4, diff_parts={K1 S1 I1 C1 F1}, diff_note={concepts: work by a battery | steps: charge on the capacitor; work equals QV}, group=TNAS02-P1-PHY-S3-G3-4]
On closing switch $S_1$ with $S_2$ open work done by battery A is (in \underline{\hspace{1.2em}} $\mu$J)
\end{bankq}
\begin{bankq}[id=TNAS02-P1-PHY-S3-Q04, source={TNAS - Test 02 (07 May 2023)}, paper=1, qno=4, type=numerical, subtopic={Energy}, answer={21.33}, answer_src=solved, working={S2 closed: C1 and C2 in series ($4/3\ \mu$F) across 4 V, $Q=16/3\ \mu$C; work by battery $=QV=64/3\approx21.33\ \mu$J}, difficulty=5, diff_parts={K1 S2 I1 C1 F1}, diff_note={concepts: capacitors in series with a battery | steps: equivalent capacitance; charge; work}, group=TNAS02-P1-PHY-S3-G3-4]
On closing switch $S_2$ with $S_1$ open work done by battery B is (in $\mu$J)
\end{bankq}
\begin{bankq}[id=TNAS02-P2-PHY-S1-Q01, source={TNAS - Test 02 (07 May 2023)}, paper=2, qno=1, type=multi, subtopic={RC circuits}, answer={C}, answer_src=key, difficulty=5, diff_parts={K1 S2 I1 C1 F1}, diff_note={concepts: RC circuit graphs | steps: time constant; initial value; compare}]
In the given circuit switch is closed at $t=0$ and capacitor is initially uncharged. Graph between $V_B-V_A$ and time is shown for different values of variable resistance $R$. Which are correct options.
\bfig{108.3}{TNAS02_p2_phy_s1_q01_1.png}
\bfig{108.3}{TNAS02_p2_phy_s1_q01_2.png}
\begin{qoptions}
\opt{A}{$|V_1|>|V_2|$}
\opt{B}{$|V_1|<|V_2|$}
\opt{C}{$R_1>R_2$}
\opt{D}{$R_1<R_2$}
\end{qoptions}
\end{bankq}
\begin{bankq}[id=TNAS02-P2-PHY-S1-Q05, source={TNAS - Test 02 (07 May 2023)}, paper=2, qno=5, type=multi, subtopic={Dielectrics}, answer={A,D}, answer_src=key, difficulty=6, diff_parts={K1 S2 I2 C1 F1}, diff_note={concepts: capacitor with partial dielectric | steps: field from the potential graph; dielectric constant; capacitance}]
A parallel plate capacitor is partially filled with dielectric. Its potential with respect to position graph is as shown. Area of plates is $2\,\mathrm m^2$.
\bfig{121.5}{TNAS02_p2_phy_s1_q05_1.png}
\begin{qoptions}
\opt{A}{Dielectric constant is $2$.}
\opt{B}{Dielectric constant is $3$.}
\opt{C}{Capacitance is $\dfrac{2000\varepsilon_0}{9}$.}
\opt{D}{Capacitance is $\dfrac{4000\varepsilon_0}{9}$.}
\end{qoptions}
\end{bankq}
\begin{bankpassage}[id=TNAS02-P2-PHY-S3-G3-4, source={TNAS - Test 02 (07 May 2023)}, questions={3-4}]
Consider a simple RC circuit as shown in figure and capacitor is initially at steady state. If by some means capacitance is suddenly reduced to one third of its initial capacitance. Initially capacitance is $2\mu$F.
\bfig{66.3}{TNAS02_p2_phy_s3_para3_4_1.png}
\end{bankpassage}
\begin{bankq}[id=TNAS02-P2-PHY-S3-Q03, source={TNAS - Test 02 (07 May 2023)}, paper=2, qno=3, type=single, subtopic={RC circuits}, answer={60}, answer_src=solved, working={Charge $60\ \mu$C is conserved while $C\to2/3\ \mu$F, so the capacitor voltage is 90 V; current $=(90-30)/1=60$ A}, difficulty=4, diff_parts={K1 S2 I1 C1 F0}, diff_note={concepts: RC circuit after a sudden change | steps: charge conserved; capacitor voltage; current}, group=TNAS02-P2-PHY-S3-G3-4]
Find current (A) in circuit just after capacitance is reduced.
\end{bankq}
\begin{bankq}[id=TNAS02-P2-PHY-S3-Q04, source={TNAS - Test 02 (07 May 2023)}, paper=2, qno=4, type=single, subtopic={RC circuits}, answer={1200}, answer_src=solved, working={Final charge $=\tfrac23\times30=20\ \mu$C, change $-40\ \mu$C; battery work $=30\times40=1200\ \mu$J}, difficulty=4, diff_parts={K1 S2 I1 C1 F0}, diff_note={concepts: battery work in a charging process | steps: final charge; charge change; work}, group=TNAS02-P2-PHY-S3-G3-4]
Find magnitude of work done by battery till steady state is achieved again in $\mu$J.
\end{bankq}
\begin{bankq}[id=TNAS03-P1-PHY-S1-Q08, source={TNAS - Test 03 (28 May 2023)}, paper=1, qno=8, type=single, subtopic={RC circuits}, answer={C}, answer_src=key, difficulty=5, diff_parts={K1 S2 I2 C1 F0}, diff_note={concepts: element in series with a resistor | steps: behaviour at early and late times; voltage graph; maximum current}]
A mysterious circuit element and a $5\ \Omega$ resistor are connected in series with a source of emf of negligible internal resistance through a switch. The element behaves like an open switch initially but allows constant current after long time. After closing the switch, a graph of voltage across the element versus time is drawn as shown in the Figure. The maximum current through the element then is :-
\bfig{106.3}{TNAS03_p1_phy_s1_q08_1.png}
\begin{qoptions}
\opt{A}{$5$A}
\opt{B}{$2$A}
\opt{C}{$3$A}
\opt{D}{$4$A}
\end{qoptions}
\end{bankq}
\begin{bankq}[id=TNAS03-P1-PHY-S1-Q12, source={TNAS - Test 03 (28 May 2023)}, paper=1, qno=12, type=single, subtopic={Charge sharing}, answer={C}, answer_src=key, difficulty=5, diff_parts={K1 S2 I2 C1 F0}, diff_note={concepts: work by a battery connected to charged plates | steps: initial and final charge; charge through the battery; work}]
Charge $Q_0$ and $2Q_0$ are given to parallel plate A and B respectively and they are separated by a small distance. The capacitance of the given arrangement is $C$. Now plates A and B are connected to positive and negative terminals of battery of potential difference $V=\dfrac{2Q_0}{C}$ respectively as shown, then the work done by the battery is :-
\bfig{63.7}{TNAS03_p1_phy_s1_q12_1.png}
\begin{qoptions}
\opt{A}{$\dfrac{2Q_0^2}{C}$}
\opt{B}{$\dfrac{4Q_0^2}{C}$}
\opt{C}{$\dfrac{5Q_0^2}{C}$}
\opt{D}{$\dfrac{6Q_0^2}{C}$}
\end{qoptions}
\end{bankq}
\begin{bankq}[id=TNAS03-P1-PHY-S1-Q13, source={TNAS - Test 03 (28 May 2023)}, paper=1, qno=13, type=single, subtopic={Combinations}, answer={C}, answer_src=key, difficulty=3, diff_parts={K1 S1 I1 C1 F0}, diff_note={concepts: capacitor network | steps: identify the symmetry; equivalent}]
Four equal capacitors, each with capacitance $C$ are connected as shown in figure. Then the equivalent capacitance between A and B is
\bfig{111.4}{TNAS03_p1_phy_s1_q13_1.png}
\begin{qoptions}
\opt{A}{$C$}
\opt{B}{$3C$}
\opt{C}{$2C$}
\opt{D}{$\dfrac{3C}{2}$}
\end{qoptions}
\end{bankq}
\begin{bankq}[id=TNAS03-P1-PHY-S2-Q07, source={TNAS - Test 03 (28 May 2023)}, paper=1, qno=7, type=numerical, subtopic={Combinations}, answer={4}, answer_src=key, difficulty=5, diff_parts={K1 S2 I1 C1 F1}, diff_note={concepts: capacitors in steady state | steps: resistor drops; voltage across the capacitor; charge}]
In the circuit shown, the charge ($\mu$C) in the 3 uF capacitor at steady state will be :-
\bfig{75.5}{TNAS03_p1_phy_s2_q07_1.png}
\end{bankq}
\begin{bankq}[id=TNAS03-P1-PHY-S2-Q08, source={TNAS - Test 03 (28 May 2023)}, paper=1, qno=8, type=numerical, subtopic={Dielectrics}, answer={16}, answer_src=key, difficulty=6, diff_parts={K1 S2 I2 C1 F1}, diff_note={concepts: capacitance with a rotating dielectric slab | steps: geometry at the angle; two regions; capacitance in parallel; evaluate}]
A cell of emf $V$ is connected to a capacitor having square plates of dimension $a$ and plate separation $d$. A dielectric slab of dielectric constant $(\sqrt3+1)$ and thickness $d$ is placed between the plates such that it exactly fills up the space between the plates at $t=0$. Now the slab is rotated with angular velocity $\omega$ as shown in the figure. Neglect edge effects the capacitance at $t=\dfrac{\pi}{6\omega}$ is $\dfrac{x\varepsilon.a^2}{8d}$. Find $x$.
\bfig{93.9}{TNAS03_p1_phy_s2_q08_1.png}
\end{bankq}
\begin{bankq}[id=TNAS03-P2-PHY-S1-Q01, source={TNAS - Test 03 (28 May 2023)}, paper=2, qno=1, type=multi, subtopic={RC circuits}, answer={B,C}, answer_src=key, difficulty=6, diff_parts={K1 S2 I2 C1 F1}, diff_note={concepts: fusing of a resistor with reversed polarity | steps: steady state voltage; power at reversal; two limits; bounds}]
A circuit diagram of an electronics device's part is given in the figure. The diagram consists of a 9 V source, a filter made of a capacitor and a resistor, a switch and output leads. The two possible positions of the switch are "on (A and B are connected) and "off" (A and C are connected). If in the given circuit the switch is turned from the position "on" to "off" (that means to turn the device off) and next the voltage source's polarity is reversed then the electronics device will still work after turning it on. However, if the voltage source's polarity is changed without turning the device off, then the resistor $R$ will fuse. Assuming that the resistor will fuse when the power on it exceeds $0.25$ W, choose the correct option(s) :-
\bfig{171.4}{TNAS03_p2_phy_s1_q01_1.png}
\begin{qoptions}
\opt{A}{The minimum value of $R$ is $400\Omega$}
\opt{B}{The minimum value of $R$ is $324\Omega$}
\opt{C}{The maximum value of $R$ is $1296\Omega$}
\opt{D}{The minimum value of $R$ is $1600\Omega$}
\end{qoptions}
\end{bankq}
\begin{bankq}[id=TNAS03-P2-PHY-S1-Q06, source={TNAS - Test 03 (28 May 2023)}, paper=2, qno=6, type=multi, subtopic={RC circuits}, answer={A,B,C}, answer_src=key, difficulty=7, diff_parts={K1 S2 I2 C2 F1}, diff_note={concepts: transient in a network with a capacitor | steps: time constant; initial values; currents at the instant; potentials}]
Five equal resistances, a capacitor and two batteries are connected in a circuit as shown in the figure. The capacitor is initially uncharged and each battery has negligible internal resistance. It is given that $R=1\Omega$, $C=1\ \mu$F, $C=1\ \mu$F and $V=8$ volt. On the basis of above information answer the following questions, if switch S is closed at $t=0$.
\bfig{103.5}{TNAS03_p2_phy_s1_q06_1.png}
\begin{qoptions}
\opt{A}{Current in the branch of capacitor at $t=5/8\ \mu$s is (approximately) 6A.}
\opt{B}{Current in the branch AB at $t=5/8\ \mu$s is (approximately) 1.7A.}
\opt{C}{Potential difference between points A and C at $t=5/8\ \mu$s is (approximately) 2.5 V}
\opt{D}{Potential difference between points A and C at $t=5/8\ \mu$s is (approximately) 6V.}
\end{qoptions}
\end{bankq}
\begin{bankq}[id=TNAS03-P2-PHY-S2-Q08, source={TNAS - Test 03 (28 May 2023)}, paper=2, qno=8, type=single, subtopic={Combinations}, answer={D}, answer_src=key, difficulty=5, diff_parts={K1 S2 I2 C1 F0}, diff_note={concepts: sequence of operations on a capacitor | steps: charge or voltage after each step; energy; match}]
\textbf{Answer the following by appropriately matching the lists based on the information given in the paragraph.}\\
A capacitor has a capacitance of 4 $\mu$F. It is filled with air and has an inter plate distance at $d$. We have a dielectric slab ($k=2$) of same thickness as the distance between the plates of the capacitor but whose area is half that of the capacitor. We perform four actions on the capacitor (i) Fill the capacitor with the dielectric, (ii) remove the dielectric slab (iii) Connect the capacitor to a battery of emf 4 V. (iv) Disconnect the battery. In column I, we mention the order of the actions. In column II, we list the various parameters after the actions. Match the entries.
\begin{center}
{\footnotesize\setlength{\tabcolsep}{2pt}\begin{tabular}{|c|c|c|p{0.4\linewidth}|}\hline
\multicolumn{2}{|c|}{List-I} & \multicolumn{2}{c|}{List-II}\\\hline
(P) & (i)(iii)(ii)(iv) & (1) & Final Charge on capacitor $=16\mu$C\\\hline
(Q) & (iii)(iv)(i)(ii) & (2) & Final energy stored in the capacitor $=32\mu$J\\\hline
(R) & (i)(iii)(iv)(ii) & (3) & Final voltage across the capacitor $=6$V\\\hline
(S) & (iii)(i)(iv)(ii) & (4) & Final Charge on capacitor $=24\ \mu$C\\\hline
 & & (5) & Final energy stored in the capacitor $=72\mu$J\\\hline
 & & (6) & Final voltage across the capacitor is 4V.\\\hline
\end{tabular}}
\end{center}
\begin{qoptions}
\opt{A}{$\mathrm P\to1,5,6;Q\to3,4,5;R\to3,4,5;S\to1,2,6$}
\opt{B}{$\mathrm P\to1,2,6;Q\to1,2,3;R\to4,5,6;S\to4,5,6$}
\opt{C}{$\mathrm P\to1,2,3;Q\to1,2,6;R\to1,5,6;S\to3,4,5$}
\opt{D}{$\mathrm P\to1,2,6;Q\to1,2,6;R\to3,4,5;S\to3,4,5$}
\end{qoptions}
\end{bankq}
\begin{bankq}[id=TNAS04-P1-PHY-S1-Q03, source={TNAS - Test 04 (18 Jun 2023)}, paper=1, qno=3, type=single, subtopic={Dielectrics}, answer={B}, answer_src=key, difficulty=5, diff_parts={K2 S2 I1 C1 F0}, diff_note={concepts: capacitance of a leaking dielectric filled capacitor, Torricelli outflow | steps: height versus time; parallel combination of filled and empty parts; substitute values; match polynomial}]
A hollow cube of side $L$ has two parallel sides conducting and remaining four sides are non conducting. It is filled with a liquid of dielectric constant $k$. A small hole of area $A$ at the bottom of the cube is opened at $t=0$ and liquid starts leaking through it. The two conducting sides are connected by a battery of emf $\varepsilon$. [$L=9$ m, $A=\dfrac{18}{\sqrt{20}}\,\mathrm{m}^2$, $k=10$, $\varepsilon=81$ v]. Find the capacitance of capacitor as a function of time ?
\bfig{87.4}{TNAS04_p1_phy_s1_q03_1.png}
\begin{qoptions}
\opt{A}{$\dfrac{\varepsilon_0}{9}\left(t^2-54t+64\right)$}
\opt{B}{$\dfrac{\varepsilon_0}{9}\left(t^2-54t+810\right)$}
\opt{C}{$\dfrac{\varepsilon_0}{9}\left(t^2+54t+810\right)$}
\opt{D}{$\dfrac{\varepsilon_0}{9}\left(t^2+54t+64\right)$}
\end{qoptions}
\end{bankq}
\begin{bankq}[id=TNAS04-P2-PHY-S1-Q04, source={TNAS - Test 04 (18 Jun 2023)}, paper=2, qno=4, type=multi, subtopic={RC circuits}, answer={A}, answer_src=key, difficulty=6, diff_parts={K1 S3 I1 C1 F1}, diff_note={concepts: capacitors in a circuit with a voltmeter and an ammeter | steps: steady state potential; reading of the voltmeter; charge through the ammeter; stored energy; compare}]
In the circuit shown in the figure, the key K is open and all capacitors are uncharged. At $t=0$, key is closed :-
\bfig{108.2}{TNAS04_p2_phy_s1_q04_1.png}
\begin{qoptions}
\opt{A}{Reading of voltmeter is $\dfrac35U_0$.}
\opt{B}{Charge flown through the ammeter is $\dfrac35CU_0$.}
\opt{C}{Energy stored in all the capacitor is $\dfrac75CU_0^2$}
\opt{D}{None of the above}
\end{qoptions}
\end{bankq}
\begin{bankq}[id=TNAS04-P2-PHY-S1-Q07, source={TNAS - Test 04 (18 Jun 2023)}, paper=2, qno=7, type=multi, subtopic={Combinations}, answer={B,C,D}, answer_src=key, difficulty=5, diff_parts={K1 S2 I1 C1 F1}, diff_note={concepts: capacitors with resistors in steady state | steps: steady state with the switch open; charge on the capacitor; after closing; current}]
In the shown circuit initially, the switch is open and the circuit is in steady state then :-
\bfig{89.6}{TNAS04_p2_phy_s1_q07_1.png}
\begin{qoptions}
\opt{A}{The charge on $3\mu$F capacitor in steady state before closing the switch is 20 $\mu$C.}
\opt{B}{The current through the battery in steady state before closing the switch is zero.}
\opt{C}{The charge on 2 $\mu$F in steady state after closing the switch is 8 $\mu$C.}
\opt{D}{The current through the battery in steady state after closing the switch is 2A.}
\end{qoptions}
\end{bankq}
\begin{bankq}[id=TNAS05-P1-PHY-S2-Q06, source={TNAS - Test 05 (09 Jul 2023)}, paper=1, qno=6, type=numerical, subtopic={Combinations}, answer={28}, answer_src=key, difficulty=7, diff_parts={K1 S3 I2 C1 F1}, diff_note={concepts: equivalent capacitance of a bridge network | steps: assign potentials; charge conservation at node one; at node two; solve potentials; charge from A; equivalent capacitance}]
If the equivalent capacitance between $A$ and $B$ (in $\mu$F) is $x$, then find the value of $10x$.
\bfig{137.6}{TNAS05_p1_phy_s2_q06_1.png}
\end{bankq}
\begin{bankq}[id=TNAS05-P1-PHY-S2-Q09, source={TNAS - Test 05 (09 Jul 2023)}, paper=1, qno=9, type=numerical, subtopic={RC circuits}, answer={24}, answer_src=key, difficulty=5, diff_parts={K1 S2 I1 C1 F1}, diff_note={concepts: time constant of an RC network | steps: capacitor replaced by open circuit; equivalent resistance; time constant; evaluate}]
In the circuit shown battery is ideal. The time constant of the circuit is $x$ C. Find the value of $x$.
\bfig{129.3}{TNAS05_p1_phy_s2_q09_1.png}
\end{bankq}
\begin{bankq}[id=TNAS05-P1-PHY-S2-Q10, source={TNAS - Test 05 (09 Jul 2023)}, paper=1, qno=10, type=numerical, subtopic={Energy}, answer={600}, answer_src=key, difficulty=6, diff_parts={K1 S2 I1 C2 F1}, diff_note={concepts: work done by a battery in a capacitor circuit | steps: charge on capacitors before; charge after; flow through the battery; work}]
Initially switch $S$ is open at $t=0$. If switch is closed then the total work done by battery 2 will be equal to \_\_\_ $\mu$J.
\bfig{126.6}{TNAS05_p1_phy_s2_q10_1.png}
\end{bankq}
\begin{bankq}[id=TNAS05-P2-PHY-S3-Q02, source={TNAS - Test 05 (09 Jul 2023)}, paper=2, qno=2, type=numerical, subtopic={Charge sharing}, answer={8}, answer_src=key, difficulty=8, diff_parts={K1 S3 I2 C2 F1}, diff_note={concepts: sequential charging of capacitors and sharing | steps: charge after the first closing; each subsequent closing; the eighth capacitor charge; capacitance with a dielectric; share charge; steady state charge; evaluate}]
In the circuit shown (figure-1), all the switches are open initially and all capacitors are uncharged. Now first only $S_1$ is closed and then opened when steady state is reached. Then $S_2$ is closed and again opened when steady state is reached, after this $S_3$ is closed and opened when again steady state is reached, this process repeated countinuously. After closing sixth switch, the fourth capacitor ($8C$) is removed and connected (say at $t=0$) at an uncharged capacitor of original capacitance $C$ which now has a dielectric (dielectric constant 2) filling half its volume (see figure-2), through a resistor $R$. If final charge on capacitor $B$ when steady state reached in figure-2 is $\dfrac n{19}CV$ . The value of $n$ is
\bfig{234.9}{TNAS05_p2_phy_s3_q02_1.png}
\bfig{92.5}{TNAS05_p2_phy_s3_q02_2.png}
\end{bankq}
\begin{bankq}[id=TNAS06-P2-PHY-S4-Q09, source={TNAS - Test 06 (30 Jul 2023)}, paper=2, qno=9, type=numerical, subtopic={Capacitors}, answer={2}, answer_src=key, difficulty=6, diff_parts={K1 S2 I2 C1 F1}, diff_note={concepts: current from a ball bouncing between charged plates | steps: charge transferred per hit; time between hits; average current; evaluate}]
Figure shows a conducting ball of capacitance $2\mu$F colliding with two large parallel plates elastically. The equivalent current in mA is :- (Assume ball moves with constant speed as shown in figure)
\bfig{160.9}{TNAS06_p2_phy_s4_q09_1.png}
\end{bankq}
\begin{bankq}[id=TNAS07-P2-PHY-S3-Q06, source={TNAS - Test 07 (20 Aug 2023)}, paper=2, qno=6, type=numerical, subtopic={Combinations}, answer={6.67}, answer_src=key, difficulty=6, diff_parts={K1 S2 I2 C1 F1}, diff_note={concepts: equivalent capacitance of a bridge network | steps: reduce series and parallel parts; symmetry of the bridge; equivalent}]
In the given network all the capacitors have same capacitance 1micro farad. Find equivalent capacitance between A and B in micro farad.
\bfig{82.3}{TNAS07_p2_phy_s3_q06_1.png}
\end{bankq}

\section*{Star Review (Other) papers}
\begin{bankq}[id=STAR03-P1-PHY-S1-Q01, source={Star Review (Other) - Test 03 (04 Dec 2023)}, paper=1, qno=1, type=single, subtopic={Capacitors}, answer={B}, answer_src=key, difficulty=3, diff_parts={K1 S2 I1 C0 F0}, diff_note={concepts: balanced bridge with capacitors | steps: condition for zero galvanometer current; potentials in terms of capacitances; ratio}]
In wheat stone's bridge shown in the adjoining figure galvanometer gives no deflection on pressing the key, the balance condition for the bridge is :-\par
\bfig{113.4}{STAR03_p1_phy_s1_q01_1.png}
\begin{qoptions}
\opt{A}{$\dfrac{R_1}{R_2}=\dfrac{C_1}{C_2}$}
\opt{B}{$\dfrac{R_1}{R_2}=\dfrac{C_2}{C_1}$}
\opt{C}{$\dfrac{R_1}{R_1+R_2}=\dfrac{C_2}{C_1-C_2}$}
\opt{D}{$\dfrac{R_1}{R_1-R_2}=\dfrac{C_2}{C_1+C_2}$}
\end{qoptions}
\end{bankq}
\begin{bankpassage}[id=STAR03-P2-PHY-S3-G3-4, source={Star Review (Other) - Test 03 (04 Dec 2023)}, questions={3-4}]
At $t=0$, switch is closed. Capacitors are initially uncharged. Current through battery at $t=0^+$ is $(10\alpha)$ Ampere while that at $t\to\infty$ is $(10/\beta)$ Ampere.
\bfig{148.9}{STAR03_p2_phy_s3_para3_4_1.png}
\end{bankpassage}
\begin{bankq}[id=STAR03-P2-PHY-S3-Q03, source={Star Review (Other) - Test 03 (04 Dec 2023)}, paper=2, qno=3, type=numerical, subtopic={RC circuits}, answer={1.50}, answer_src=key, difficulty=7, diff_parts={K1 S2 I2 C2 F1}, diff_note={concepts: initial current in a capacitor network | steps: capacitors as short circuits at the start; equivalent resistance; current}, group=STAR03-P2-PHY-S3-G3-4]
The value of $\alpha$ is \underline{\hspace{1cm}}
\end{bankq}
\begin{bankq}[id=STAR03-P2-PHY-S3-Q04, source={Star Review (Other) - Test 03 (04 Dec 2023)}, paper=2, qno=4, type=numerical, subtopic={RC circuits}, answer={3.00}, answer_src=key, difficulty=7, diff_parts={K1 S2 I2 C2 F1}, diff_note={concepts: steady current in a capacitor network | steps: capacitors as open circuits; path through the resistors; current; value of beta}, group=STAR03-P2-PHY-S3-G3-4]
Value of $\beta$ is \underline{\hspace{1cm}}
\end{bankq}

\section*{Internal Test papers}
\begin{bankq}[id=INT02-P1-PHY-S3-Q16, source={Internal Test - Test 02 (07 Jul 2021)}, paper=1, qno=16, type=numerical, subtopic={Charge sharing}, answer={0.83}, answer_src=key, difficulty=7, diff_parts={K1 S2 I2 C2 F1}, diff_note={concepts: connected charged spheres with equal surface charge density | steps: charges of both; common potential; charge sharing ratio; new surface density}]
Two isolated metal spheres of radii $R$ and $2R$ are charged such that both have the same surface charge density $\sigma$. The spheres are located far away from each other. When they are connected by thin conducting wire, the new surface charge density on the bigger sphere will be $a\sigma$ then find $a$.
\end{bankq}
\begin{bankq}[id=INT03-P1-PHY-S1-Q01, source={Internal Test - Test 03 (08 Sep 2021)}, paper=1, qno=1, type=single, subtopic={RC circuits}, answer={B}, answer_src=key, difficulty=7, diff_parts={K1 S3 I2 C2 F0}, diff_note={concepts: charging and discharging cycle of a capacitor between two thresholds | steps: charging law between the two voltages; discharging law; time of each part; sum; period}]
The switch in circuit shifts from 1 to 2 when $V_C>2V/3$ and goes back to 1 from 2 when $V_C<V/3$. The voltmeter reads voltage as plotted. What is the period $T$ of the wave form in terms of $R$ and $C$?\par
\bfig{278.8}{INT03_p1_phy_s1_q01_1.png}

\begin{qoptions}
\opt{A}{$RC\ln 3$}
\opt{B}{$2RC\ln 2$}
\opt{C}{$\dfrac{RC}{2}\ln 3$}
\opt{D}{$\dfrac{RC}{3}\ln 3$}
\end{qoptions}
\end{bankq}
\begin{bankq}[id=INT03-P1-PHY-S2-Q05, source={Internal Test - Test 03 (08 Sep 2021)}, paper=1, qno=5, type=multi, subtopic={Capacitors}, answer={A}, answer_src=key, difficulty=5, diff_parts={K1 S2 I2 C0 F1}, diff_note={concepts: capacitor with plates not parallel | steps: potential difference fixed; field at each point; surface charge density; compare}]
In the given figure, a capacitor of non-parallel plates is shown. The plates of capacitor are connected by a cell of emf $V_0$. If $\sigma$ denotes surface charge density and $E$ denotes electric field. Then\par
\bfig{111.3}{INT03_p1_phy_s1_q05_1.png}
\begin{qoptions}
\opt{A}{$\sigma_A>\sigma_B$}
\opt{B}{$E_F>E_D$}
\opt{C}{$E_F=E_D$}
\opt{D}{$\sigma_A=\sigma_B$}
\end{qoptions}
\end{bankq}
\begin{bankq}[id=INT03-P1-PHY-S2-Q06, source={Internal Test - Test 03 (08 Sep 2021)}, paper=1, qno=6, type=multi, subtopic={Charge sharing}, answer={A,B,C}, answer_src=key, difficulty=8, diff_parts={K1 S3 I2 C2 F1}, diff_note={concepts: charge redistribution between capacitors and a third capacitor | steps: charges of the first two; connection with opposite polarity; common potential; third capacitor; direction of flow}]
Two capacitors of $2\,\mu$F and $3\,\mu$F are charged to 150 volt and 120 volt respectively. The plates of capacitor are connected as shown in the figure. A discharged capacitor of capacity $1.5\,\mu$F falls to the free ends of the wire. Then --\par
\bfig{126.1}{INT03_p1_phy_s1_q06_1.png}
\begin{qoptions}
\opt{A}{Charge on the $1.5\,\mu$F capacitor is $180\,\mu$C}
\opt{B}{Charge on the $2\,\mu$F capacitor is $120\,\mu$F}
\opt{C}{Charge flows through A from right to left}
\opt{D}{Charge flows through A from left to right.}
\end{qoptions}
\end{bankq}
\begin{bankq}[id=INT03-P1-PHY-S3-Q13, source={Internal Test - Test 03 (08 Sep 2021)}, paper=1, qno=13, type=numerical, subtopic={Dielectrics}, answer={9}, answer_src=key, difficulty=6, diff_parts={K1 S2 I2 C1 F1}, diff_note={concepts: capacitance of two dielectric regions | steps: geometry of the regions; capacitance of each; ratio}]
Two medium of dielectric constant $K_1$ and $K_2$ are introduced according to given figure. If $\dfrac{K_1}{K_2}=3$ then calculate ratio of capacity of part A and part B. (Area of each plate is $A$)\par
\bfig{89.3}{INT03_p1_phy_s2_q13_1.png}
\end{bankq}

\section*{Nurture (General) papers}
\begin{bankq}[id=NUR01-P1-PHY-S2-Q01, source={Nurture (General) - Test 01 (12 Jun 2021)}, paper=1, qno=1, type=single, subtopic={RC circuits}, answer={B}, answer_src=solved, working={Compare the initial currents in the four circuits}, difficulty=6, diff_parts={K2 S2 I2 C1 F0}, diff_note={concepts: initial currents in RC style circuits, rc circuits | steps: identify the elements of each case; initial current; compare}]
Of the given 4 situations, which one has maximum value of current just after closing the switch ?
\begin{qoptions}
\opt{A}{I, iii, P}
\opt{B}{IV, i, Q}
\opt{C}{II, iii, P}
\opt{D}{III, i, S}
\end{qoptions}
\end{bankq}
\begin{bankq}[id=NUR01-P1-PHY-S2-Q02, source={Nurture (General) - Test 01 (12 Jun 2021)}, paper=1, qno=2, type=single, subtopic={RC circuits}, answer={A}, answer_src=solved, working={Charge drawn in one time constant}, difficulty=6, diff_parts={K2 S2 I2 C1 F0}, diff_note={concepts: charge drawn in one time constant, rc circuits | steps: identify each case; charge function; compare}]
In which case the charge drawn from the battery(s) during $t=0$ to $t=RC$ is the least ?
\begin{qoptions}
\opt{A}{II, i, S}
\opt{B}{I, iii, P}
\opt{C}{III, ii, Q}
\opt{D}{IV, i, Q}
\end{qoptions}
\end{bankq}
\begin{bankq}[id=NUR01-P1-PHY-S2-Q03, source={Nurture (General) - Test 01 (12 Jun 2021)}, paper=1, qno=3, type=single, subtopic={RC circuits}, answer={D}, answer_src=solved, working={Current differs between t=0 and long times only for the capacitor circuit}, difficulty=6, diff_parts={K2 S2 I2 C1 F0}, diff_note={concepts: current at start and after a long time, rc circuits | steps: initial current; final current; compare}]
In which case the current through battery(s) has different non zero values just after closing the switch \& after a long time ?
\begin{qoptions}
\opt{A}{II, i, S}
\opt{B}{I, iv, Q}
\opt{C}{IV, i, Q}
\opt{D}{III, iii, R}
\end{qoptions}
\end{bankq}
\begin{bankq}[id=NUR01-P2-PHY-S2-Q02, source={Nurture (General) - Test 01 (12 Jun 2021)}, paper=2, qno=2, type=multi, subtopic={RC circuits}, answer={A,B,C}, answer_src=solved, working={Charges equalise to (Q1+Q2)/2; initial potential difference (Q1-Q2)/2C}, difficulty=7, diff_parts={K2 S2 I2 C1 F1}, diff_note={concepts: capacitor with unequal plate charges, rc circuits | steps: equalisation of charges; time dependence; initial potential difference}]
The two plates of an ideal capacitor of capacitance $C$ are given charges $Q_1$ and $Q_2$ as shown in the circuit below. This capacitor is connected across a resistance $R$ as shown key is closed at $t=0$. Then after time $t$.\par
\bfig{66.1}{NUR01_p2_phy_s2_q02_1.png}
\begin{qoptions}
\opt{A}{Total charge on the right plate is $q_2=\dfrac{Q_1+Q_2}2-\left(\dfrac{Q_1-Q_2}2\right)e^{-\frac t{RC}}$}
\opt{B}{Total charge on the left plate is $q=\dfrac{Q_1+Q_2}2+\left(\dfrac{Q_1-Q_2}2\right)e^{-\frac t{RC}}$}
\opt{C}{Initial potential difference across the plates is given by $\left|\dfrac{Q_1-Q_2}{2C}\right|$}
\opt{D}{Initial potential difference across the plates is given by $\dfrac{Q_1+Q_2}{2C}$}
\end{qoptions}
\end{bankq}
\begin{bankq}[id=NUR03-P1-PHY-S1-Q15, source={Nurture (General) - Test 03 (27 Sep 2023)}, paper=1, qno=15, type=single, subtopic={RC circuits}, answer={C}, answer_src=key, difficulty=5, diff_parts={K2 S2 I1 C1 F0}, diff_note={concepts: charging with only one switch closed, rc circuits | steps: charge growth law; energy accounting; check the options}]
In an LCR circuit as shown below both switches are open initially. Now switch $S_1$ is closed, $S_2$ kept open, ($q$ is charge on the capacitor and $\tau=RC$ is Capacitive time constant). Which of the following statement is correct?\par
\bfig{92.1}{NUR03_p1_phy_s1_q15_1.png}
\begin{qoptions}
\opt{A}{Work done by the battery is half of the energy dissipated in the resistor}
\opt{B}{At $t=\tau$, $q=CV/2$}
\opt{C}{At $t=2\tau$, $q=CV\left(1-e^{-2}\right)$}
\opt{D}{At $t=2\tau$, $q=CV\left(1-e^{-1}\right)$}
\end{qoptions}
\end{bankq}
\end{document}
